% !PRINTED 531
\weberpage{531}
The discriminant $\varDelta(\omega_1', \omega_2', \ldots, \omega_n')$ is therefore smaller than $\varDelta(\omega_1, \omega_2, \ldots, \omega_n)$, and this contradicts the assumption that $\omega_1, \omega_2, \ldots, \omega_n$ is a minimal basis. With this our theorem is proved.

If one denotes the totality of all \gesperrt{integers} of the field $\Omega$ by $\mathfrak{o}$, then one obtains all the numbers of $\mathfrak{o}$, and each only once, if in
\weq{6}{k_1\omega_1 + k_2\omega_2 + \cdots + k_n\omega_n}
one lets the coefficients $k_1, k_2, \ldots, k_n$ run through all the integral rational numbers.

For this reason a minimal basis of $\Omega$ is also called a \gesperrt{basis of $\mathfrak{o}$}.

\gesperrt{The discriminant of a basis of $\mathfrak{o}$ is equal to the fundamental number of the field $\Omega$.}

By formula (11) of the preceding section we can derive from a basis of $\mathfrak{o}$ arbitrarily many others by linear substitution:
\weq{7}{(\omega_1', \omega_2', \ldots, \omega_n') = C(\omega_1, \omega_2, \ldots, \omega_n),}
if $C$ (in the notation of \S.\,37) denotes a linear substitution with rational integral coefficients and the determinant $\pm 1$.

Since, by the meaning of the basis of $\mathfrak{o}$, all integers in $\Omega$, and hence also the elements $\omega_1', \omega_2', \ldots, \omega_n'$ of a second basis of $\mathfrak{o}$, are representable linearly with integral rational coefficients by $\omega_1, \omega_2, \ldots, \omega_n$, it follows conversely that one can derive all the others from one basis of $\mathfrak{o}$ by such linear substitutions with the determinant $\pm 1$.

For if
\[
(\omega_1, \omega_2, \ldots, \omega_n) = C'(\omega_1', \omega_2', \ldots, \omega_n'),
\]
then the composed substitution $CC'$ must be the identical one, and the product of the two determinants is therefore 1, and hence each of them $= \pm 1$.

Since among the numbers of $\mathfrak{o}$ the number 1 is always contained, one can, if $\omega_1, \omega_2, \ldots, \omega_n$ is a basis of $\mathfrak{o}$, determine the \gesperrt{integral rational numbers} $c_1, c_2, \ldots, c_n$ so that the relation
\weq{8}{c_1\omega_1 + c_2\omega_2 + \cdots + c_n\omega_n = 1}
is satisfied.
